% Imporditud: tools/import_eksperimendid.py
% Allikas: Eesti füüsikaolümpiaadi eksperimendi ülesannete
%   kogu (Rael Kalda, 2023), ülesanne Ü119, lahendus L119.
\setAuthor{Koit Timpmann}
\setRound{piirkonnavoor}
\setYear{2020}
\setNumber{G E2}
\setDifficulty{3}
\setTopic{elektriahelad}
\setVahendid{Taskulambipirn, patarei, ühendusjuhtmed (4 tk.), multimeeter}
\setPunktid{12}
\setAllikas{Eksperimendi kogu Ü119, lahendus lk 94}
\setLahendusseis{toores}

\prob{Taskulambipirni töötemperatuur}
Määrake töötava taskulambipirni hõõgniidi temperatuur. Volframi eritakistuse temperatuurikoefitsent $\alpha$ = 0,0044 1/K. Hõõgniidi soojuspaisumist ei ole vaja arvestada.
\textit{Märkus:} Eeldage, et eritakistuse jaoks kehtib valem $\rho$ = $\rho_0$ (1 + $\alpha$t), kus t on temperatuur Celsiuse kraadides, $\rho_0$ on aine eritakistus temperatuuril 0 $^\circ$C ja $\alpha$ on eritakistuse temperatuuritegur.

\hint

\solu
Taskulambipirni hõõgniidi takistus sõltub temperatuurist. Teades toatemperatuuri, hõõgniidi takistust toatemperatuuril ja hõõgniidi aine eritakistuse temperatuurikoefitsienti on võimalik arvutada põleva taskulambipirni hõõgniidi temperatuur. (2 p) Tähistame: toatemperatuur - $t_1$, põleva lambi hõõgniidi temperatuur - $t_2$, hõõgniidi takistus toatemperatuuril $R_1$, põleva lambi hõõgniidi takistus $R_2$. Lambi hõõgniidi takistus toatemperatuuril on

$R_1$ = $\rho_0$ (1 + $\alpha$$t_1$) l/S (0,5 p) Põleva lambi hõõgniidi takistus on

$R_2$ = $\rho_0$ (1 + $\alpha$$t_2$) l/S (0,5 p) avaldame temperatuuri $t_2$:

$R_2$/$R_1$ = 1 + $\alpha$$t_2$ (1 p) 1 + $\alpha$$t_1$

$t_2$ = $R_2$ + $R_2$$\alpha$$t_1$ $-$ $R_1$ (1 p) $R_1$$\alpha$

Mõõtmised:

1. Valime oommeetri sobiva mõõtepiirkonna ja mõõdame toatemperatuuril oleva taskulambipirni takistuse $R_1$ (1 p).

2. Ühendame taskulambipirni patareiga, valime voltmeetri sobiva mõõtepiirkonna ja mõõdame pinge taskulambipirni klemmidel. (1 p)

3. Ühendame patareiga jadamisi taskulambipirni ja ampermeetri, valime ampermeetri sobiva piirkonna ja mõõdame voolutugevuse hõõgniidis. (1 p)

Arvutused:

1. Arvutame põleva lambi hõõgniidi takistuse valemi $R_2$ = U/I abil. (1 p)

2. Arvutame põleva lambi hõõgniidi temperatuuri. (1 p)

Kordusmõõtmised (1 p). Sõltuvalt taskulambipirnist on hõõgniidi töötemperatuur vahemikus 1400 $^\circ$C - 1800 $^\circ$C (1 p).
\probend
